% !TEX root = ../main.tex
\PlanPage{政策事件案例}{武桐宇}{28}
\textbf{本页结论：}政策案例的研究价值在于识别政策所指向的基础设施、场景与技术环节，而不是把政策目标直接写成公司订单、收入或拨款。

广东省政府办公厅于 2024 年 5 月发布《广东省推动低空经济高质量发展行动方案（2024—2026年）》，其文件信息显示成文日期为 2024 年 5 月 21 日、发布日期为 5 月 22 日\cite{wty-gd-plan}。方案提出起降设施、低空智联网、数据底座、应用场景与产业链协同等方向。此处使用的是政策原文的方向性表述，而非对任一项目进度的推断。

\begin{table}[htbp]
\centering
\caption{广东低空行动方案的研究映射示例}
\label{tab:wty-policy-case}
\small
\begin{tabularx}{\linewidth}{p{0.25\linewidth}X p{0.28\linewidth}}
\toprule
政策内容（原文方向） & 可进入图谱的研究节点 & 不能直接推出的事实 \\
\midrule
起降场、能源站与保障设施 & 基础设施、充换电、运营保障 & 某企业中标、开工或确认收入 \\
通信、导航、监视与数据底座 & 飞控/通信、低空智联网、数据治理 & 特定设备已部署或形成市场份额 \\
物流、客运、应急等场景 & 运营与应用节点、整机适用场景 & 航线已获批准、已常态化运营 \\
关键技术与产业链协同 & 动力、控制、整机与零部件节点 & 企业获得补贴或技术已成熟 \\
\bottomrule
\end{tabularx}
\end{table}

可复现的研究路径是：保存政策原文与适用地区，提取条款中的对象词，映射到图谱节点，再逐条搜索公司公开材料验证关联。若没有公司层面的直接证据，输出停在“候选研究对象”，不生成暴露度或投资判断。
